%% kapa: exam=1 round=27 year=2016 period=2 form=A subject=accounting subject_ko=회계학 seq=18 number=58
%% source: Q-Net 감정평가사 기출문제 (cst00309) artlSeq=5169402 file=extracted/1008/5169402/제27회 감정평가사 1차 문제/제27회 감정평가사 2교시문제(A형).pdf
%% conversion: kordoc 4.24.51 -> tools/gen1.py (auto); review=proofread-sonnet; flags=none
\begin{question}{58}
공정가치 측정에 관한 설명으로 옳지 않은 것은?
\begin{choices}
  \item 공정가치 측정은 자산을 매도하거나 부채를 이전하는 거래가 주된 시장이나 가장 유리한 시장(주된 시장이 없는 경우)에서 이루어지는 것으로 가정한다.
  \item 부채의 공정가치는 불이행위험의 효과를 반영한다.
  \item 자산이나 부채의 공정가치를 측정하기 위하여 사용되는 주된 시장의 가격에서 거래원가는 조정한다.
  \item 요구불 특성을 가진 금융부채(예: 요구불예금)의 공정가치는 요구하면 지급요구가 가능한 최초일부터 할인한 금액 이상이어야 한다.
  \item 자산이나 부채의 공정가치는 자산을 매도하면서 수취하거나 부채를 이전하면서 지급하게 될 가격(유출가격)이다.
\end{choices}
\end{question}
\answer{3}
